\documentclass[10pt]{article}
\usepackage[margin=0.82in]{geometry}
\usepackage{amsmath,amssymb,booktabs,array,microtype}
\usepackage[T1]{fontenc}
\usepackage{lmodern}
\usepackage[hidelinks]{hyperref}
\usepackage{xurl}
\usepackage{enumitem}
\setlist{nosep,leftmargin=*}
\setlength{\parindent}{0pt}
\setlength{\parskip}{4pt}
\newcommand{\Sep}{\operatorname{Sep}}
\newcommand{\score}{\operatorname{score}}

\title{\textbf{Exact Computation of Near-Identities for Five-State Binary Automata}}
\author{Jinji Li\\\small Independent Researcher, Hong Kong}
\date{September 2026}

\begin{document}
\maketitle
\vspace{-1.4em}

\begin{abstract}
We study separation of two binary words by deterministic finite automata (DFAs), equivalently identities in finite transformation semigroups.  We implement a reproducible enumeration of accessible binary DFA transition structures up to five states and evaluate word pairs against the complete canonical universe of 166,152 structures.  As a validation, the implementation independently verifies the known length-48 obstruction for five-state automata.  We then investigate length-47 pairs immediately below the proven threshold.  A bounded computational search produces a pair separated by only 52 canonical structures; all 52 require five states, one letter acts as a permutation, and the other acts as a rank-four transformation with indegree profile $(0,1,1,1,2)$.  Structured perturbations exhibit witness exchange: eliminating one exceptional separator set can create a disjoint replacement set.  A natural structural conjecture suggested by this behavior is then falsified by exact computation.  We make no claim that score 52 is globally optimal among length-47 pairs.  The contribution is an exact, reproducible computational description of a near-identity and its separator landscape, together with machine-checkable artifacts for further structural study.
\end{abstract}

\section{Introduction}
Given two distinct words $u,v$, let $\Sep(u,v)$ be the minimum number of states in a deterministic finite automaton that accepts exactly one of them.  A transition structure permits such a separation exactly when the two words terminate in different states: an accepting set can then be chosen to contain exactly one of the two endpoints.  Accepting-state subsets therefore need not be enumerated.

The problem has a direct algebraic formulation.  Let $T_k$ be the full transformation semigroup on $k$ elements.  Bulatov, Karpova, Shur, and Startsev showed that identities in $T_k$ correspond to word pairs that cannot be separated by a $k$-state DFA, and constructed a family of short identities in $T_k$ \cite{BKSS2017}.  For $k=5$ their construction gives a binary identity of length 48.  They conjectured that this identity was shortest.  Karpova and Shur subsequently proved that the shortest identity in $T_5$ has length 48 \cite{KS2021}.

Thus the extremal threshold for $T_5$ is already known.  The present work asks a narrower computational question: \emph{what does the separator landscape look like immediately below that threshold?}  We build an exact canonical enumeration of accessible binary DFA transition structures with at most five states, reproduce the known length-48 obstruction, and use this finite universe to study length-47 near-identities.  Our main object is an explicit length-47 pair with only 52 canonical separators.  We characterize these exceptional witnesses and test a structural hypothesis suggested by their behavior.

The purpose of this note is computational and structural.  In particular, we do not reprove the length-48 optimality theorem, do not claim that 52 is the minimum possible score at length 47, and do not claim a universal witness-exchange theorem.

\section{Exact canonical enumeration}
A binary DFA transition structure with state set $Q=\{0,\ldots,k-1\}$ is determined by two transformations
\[
T_0,T_1:Q\to Q,
\]
with state $0$ fixed as the initial state.  We enumerate only accessible structures.  Starting with state $0$ discovered, scan transition entries in the fixed order $(0,0),(0,1),(1,0),(1,1),\ldots$, where each entry specifies a source state and an input symbol.  If $d$ states have been discovered, the next target may be any label in $\{0,\ldots,d-1\}$ or, when $d<k$, the new label $d$.  Choosing $d$ increases the number of discovered states by one.  A branch is discarded if the next source state has not been discovered; a completed table is emitted only when all $k$ states have been discovered.

For any accessible structure, this fixed scan assigns each state its label at first discovery.  Relabeling states while preserving the initial state and the two input symbols does not change the resulting canonical table.  Conversely, identical canonical tables describe isomorphic pointed transition structures.  The procedure thus removes state-relabeling redundancy without identifying the two input symbols.  The direct generator is checked against a separate labelled-table enumeration through three states; this cross-check has not been extended to five states.  The resulting counts are shown in Table~\ref{tab:counts}.

\begin{table}[h]
\centering
\caption{Canonical accessible binary transition structures.}
\label{tab:counts}
\begin{tabular}{rr}
\toprule
States $k$ & Number of structures \\
\midrule
1 & 1 \\
2 & 12 \\
3 & 216 \\
4 & 5,248 \\
5 & 160,675 \\
\midrule
Total & 166,152 \\
\bottomrule
\end{tabular}
\end{table}

For reproducibility, the ordered five-state transition-table enumeration used in the experiments has SHA-256 fingerprint
\begin{center}
\small\ttfamily
9982601f70cbf4f4deeed5f748db76352fa8700a52d6980a990a5a4b184b5d75.
\end{center}
For a transition structure $A$, write $\delta_A(0,w)$ for the state reached from the initial state after reading $w$.  We define the canonical separator score
\[
\score(u,v)=\left|\left\{A\in\mathcal A_{\leq 5}:\delta_A(0,u)\neq\delta_A(0,v)\right\}\right|,
\]
where $\mathcal A_{\leq5}$ is the complete set of 166,152 canonical accessible structures.  This score is a computational statistic introduced in this note: it counts separating transition structures, not accepting-state subsets, and is not $\Sep(u,v)$, which minimizes the number of states.  Any DFA can be restricted to the states reachable from its initial state without changing the two runs.  A score of zero is therefore equivalent to failure of every DFA with at most five states to separate the pair.

\subsection{Regression against the known length-48 identity}
For $k=5$, the construction of \cite{BKSS2017} specializes to
\[
U=(01)^{15}(10)^5(01)^4,\qquad
V=(01)^3(10)^5(01)^{16}.
\]
Both words have length 48.  Exhaustive evaluation of $(U,V)$ against the canonical universe gives zero separators for each state size $k=1,\ldots,5$.  Hence none of the 166,152 structures separates the pair.  This independently reproduces the known obstruction and serves as a regression test for the enumeration and evaluation pipeline; the mathematical result itself is not new.

\section{An explicit length-47 near-identity}
Our bounded experiments included one-deletion neighborhoods of the length-48 construction, local and structured perturbations, witness-guided search, counterexample-guided search, and shared-block families.  These searches were not exhaustive over all length-47 word pairs.

The strongest pair found in these experiments is
\begin{align*}
u={}&\texttt{10101010101010101010101010101101010101001010101},\\
v={}&\texttt{10101101010101001010101010101010101010101010101}.
\end{align*}
Both words have length 47.  Independent exhaustive reevaluation against all 166,152 canonical structures gives
\[
\boxed{\score(u,v)=52}.
\]
The separator counts for state sizes $1,2,3,4,5$ are respectively
\[
(0,0,0,0,52).
\]
Thus $\Sep(u,v)=5$: within the enumerated universe all 52 separating structures have five states, while 166,100 of the 166,152 canonical structures fail to separate the pair.  This is $99.96870\%$ of the canonical universe.

The number 52 is a property of this explicit pair, not a claimed extremal value.  The search does not cover all ordered length-47 pairs, so no global optimality conclusion follows.

\subsection{Structure of the 52 separators}
The complete set of 52 separators exhibits a uniform transformation pattern.  For every witness,
\[
T_1\in S_5,\qquad \operatorname{rank}(T_0)=4,
\]
where $T_1$ is a permutation of the five states.  The sorted indegree profile of $T_0$ is
\[
(0,1,1,1,2),
\]
so $T_0$ has one missing image and one collision.  The observed cycle types of $T_1$ are
\[
(1,4),\ (5),\ (2,3),\ (1,1,3),\ (1,2,2),\ (1,1,1,2).
\]
These statements are an exact census of the 52 witnesses for this pair; they are not asserted for arbitrary length-47 pairs.

This pattern is reminiscent of the role of permutation and non-permutation cases in the proof of the length-48 construction \cite{BKSS2017}, but the statements concern different objects.  The earlier proof classifies the action arising in the constructed identity, whereas the census above classifies the exceptional automata that separate one particular near-identity.  We therefore do not identify the two mechanisms.

\section{Witness exchange and an exact negative result}
The near-identity admits a useful shared-block description,
\[
u=AB,\qquad v=BC,\qquad A=(10)^{12},\quad C=(01)^{12},
\]
where $B=\texttt{10101101010101001010101}$ has length 23.  Consider the restricted family $(AB',B'C)$ in which $B'$ is obtained from $B$ by flipping at most two bits.  The same modified block appears as a suffix of the first word and a prefix of the second.  Let $W_0$ denote the 52 canonical separators of the original pair.  In this family, some perturbations eliminate every member of $W_0$ yet introduce new separators; one pair has score 64 with a separator set disjoint from $W_0$.

The restricted implication tested was: \emph{for every pair $(AB',B'C)$ in this family, if no member of $W_0$ separates it, then at least one replacement separator with at most five states has a singular symbol map}, meaning that $T_0$ or $T_1$ has rank smaller than the number of states.  Here a replacement is a canonical separator outside $W_0$.

An exact counterexample is obtained by flipping bit 5 of $B$, numbering positions from zero, to give $B'=\texttt{10101001010101001010101}$.  None of the 52 original witnesses separates $(AB',B'C)$, but exhaustive evaluation gives $\score(AB',B'C)=74$.  All 74 replacement witnesses lie in the relevant permutation branch: both $T_0$ and $T_1$ permute five states, and the map induced by reading $10$, namely $T_0\circ T_1$, is a full 5-cycle.  This branch consists of 120 canonical accessible transition tables.  Hence the counterexample has no separator with a singular symbol map in the universe of at most five states, and the displayed implication is false.

This falsification concerns only the stated shared-block implication.  The original witnesses themselves have a singular $T_0$; the experiment does not establish a general implication from eliminating permutation witnesses to producing singular witnesses, nor does it rule out every weaker structural statement.

\section{Reproducibility and limitations}
The implementation and machine-readable artifacts are maintained at
\begin{center}
\url{https://github.com/JinjiLI-0725/separating-words}.
\end{center}
At the v1 research freeze, the regression suite contains 23 tests and all 23 pass.  The repository includes canonical enumeration code, the length-48 verifier, saved witness sets and summaries, structural-analysis scripts, and the exact counterexample used to falsify the restricted witness-exchange conjecture.

All reported scores count canonical accessible transition structures, with initial state fixed at $0$ and input symbols distinguished.  No multiplicity is assigned for the different accepting sets that could separate the same two endpoints.  Scalar and vectorized endpoint evaluations provide separate evaluation paths but share the direct canonical generator.

The principal limitations are deliberate.  First, the search producing the score-52 pair is bounded rather than exhaustive over all length-47 pairs.  Second, several V6 perturbation observations are empirical statements about explicitly enumerated structured families, not universal theorems.  Third, the 2021 optimality result is cited as prior work; the present computation neither rederives nor strengthens that theorem.  Finally, the structural regularity of the 52 witnesses remains an observation about one explicit near-identity.  A mathematical characterization of when such permutation/rank-deficient witness families must occur remains open for further investigation.

\section{Conclusion}
Exact enumeration provides a small but useful laboratory for the five-state word-separation problem.  It reproduces the known length-48 obstruction, identifies an explicit length-47 pair with only 52 canonical separators, and reveals a highly constrained witness structure.  Subsequent perturbation experiments show that witness sets can be exchanged rather than monotonically eliminated, and exact falsification rules out a natural simple explanation of that behavior.

The resulting picture is intentionally narrower than an extremal theorem: immediately below a proven identity threshold, near-identities can have very small, structured, and unstable exceptional separator sets.  The accompanying artifacts make these observations independently checkable and provide a concrete basis for future structural arguments.

\begin{thebibliography}{9}
\bibitem{BKSS2017}
A. A. Bulatov, O. Karpova, A. M. Shur, and K. Startsev,
``Lower Bounds on Words Separation: Are There Short Identities in Transformation Semigroups?,''
\emph{The Electronic Journal of Combinatorics}, vol. 24, no. 3, P3.35, 2017.
\href{https://doi.org/10.37236/6450}{doi:10.37236/6450}.

\bibitem{KS2021}
O. Karpova and A. M. Shur,
``Words Separation and Positive Identities in Symmetric Groups,''
\emph{Journal of Automata, Languages and Combinatorics}, vol. 26, nos. 1--2, pp. 67--89, 2021.
\href{https://doi.org/10.25596/jalc-2021-067}{doi:10.25596/jalc-2021-067}.
\end{thebibliography}

\end{document}
